\begin{textsample}{NEG Dim 2 – 2017\_c5 – Score -3.00 – 2017\_c5/enade\_2017\_fis\_eeb3bfda27.txt}  \label{ex:f2_neg_046}
CAPÍTULO 5 DISTRIBUIÇÃO DOS CONCEITOS 5.1 PANORAMA NACIONAL DA DISTRIBUIÇÃO DOS CONCEITOS O Conceito Enade é um indicador de qualidade que avalia os cursos por intermédio dos desempenhos dos estudantes a partir dos resultados do Enade.
Os valores possíveis vão de 1 ( pior situação ) a 5 ( melhor situação ).
Os cursos com apenas um ou sem nenhum concluinte participante não obtêm o Conceito Enade, ficando Sem Conceito ( SC ).
Na Tabela 5.1a, são apresentadas a quantidade e a distribuição de cursos de Licenciatura de Física participantes do Enade/2017, por faixa de conceito e Grande Região. \textbf{Enfatiza} -se, mais uma vez, que a diferença dentre os cursos tabulados neste capítulo e no capítulo 2 corresponde aos cursos Sem Conceito ( SC ), ou seja, àqueles, em princípio, sem alunos concluintes que houvessem participado da prova ou àqueles que tivessem tão somente um aluno concluinte18.
Observando -se os dados da Tabela 5.1a, nota -se que, dos 217 cursos participantes, 79 ( 36,4 \% ) classificaram -se com conceito 3, o valor modal.
Esse foi também o conceito modal nas regiões Sudeste ( 41,8 \% ) e Centro-Oeste ( 42,1 \% ).
Na região Sul, a moda foi o conceito 4, atribuído a 41,5 \% dos cursos da região.
Nas regiões Norte e Nordeste, a moda foi o conceito 2, atribuído a 69,6 \% e 44,8 \% dos cursos da região, respectivamente.
O conceito 2 foi o segundo mais frequente em nível nacional ( 29,0 \%, correspondendo a 63 cursos ) e o conceito 4, o terceiro ( 21,2 \%, correspondendo a 46 cursos ).
Houve, ainda, oito cursos ( 3,7 \% ) que receberam conceito 1, e nove cursos ( 4,1 \% ) que receberam conceito 5.
Dos 217 cursos de Física ( Licenciatura ), 12 ( 5,5 \% ) ficaram Sem Conceito ( SC ). 18 Estes cursos com somente um concluinte não foram considerados no capítulo 2.
A região Norte participou com 23 cursos ou 10,6 \% do total nacional.
Desses, 16 cursos receberam o conceito 2, o que equivale a 69,6 \% do total regional, além de ser o conceito modal na região, como já comentado.
A um curso ( 4,3 \% ) atribuiu -se conceito 4, e a três cursos cada ( 13,0 \% ) os conceitos 1 e 3.
Nenhum curso ficou Sem Conceito ( SC ) ou recebeu conceito 5 nessa região.
A região Nordeste participou com 67 cursos ou 30,9 \% do total nacional.
Nessa região, 30 cursos ( 44,8 \% em termos regionais ) obtiveram conceito 2, o conceito modal para a região, como já comentado.
Os conceitos 1 e 3 foram atribuídos, respectivamente, a dois e a 28 cursos ( 3,0 \% e 41,8 \%, respectivamente ).
Ainda nessa região, quatro cursos ( 6,0 \% ) foram avaliados com o conceito 4 e a nenhum curso foi atribuído o conceito 5.
Três cursos ( 4,5 \% ) ficaram Sem Conceito ( SC ).
Dos 67 cursos participantes da região Sudeste ( 30,9 \% do total nacional ), o conceito 3 foi o conceito modal, como já anteriormente assinalado, tendo sido obtido por 28 cursos ( 41,8 \% ).
O conceito 1 foi atribuído a um curso ( 1,5 \% ), e o conceito 2, a oito cursos ( 11,9 \% ).
Receberam os conceitos 4 e 5, respectivamente, 21 cursos ( 31,3 \% ) e seis cursos ( 9,0 \% ).
Nessa região, três cursos ( 4,5 \% ) ficaram Sem Conceito ( SC ).
Os 41 cursos da região Sul corresponderam a 18,9 \% do total nacional.
Como já indicado, houve predominância do conceito 4, atribuído a 41,5 \% dos cursos da região ( conceito modal ), ou seja, atribuído a 17 dos 41 cursos participantes na região Sul.
O conceito 2 foi atribuído a três cursos ( 7,3 \% ), o conceito 3, a 12 cursos ( 29,3 \% ), e o conceito 5, a três cursos ( 7,3 \% ).
Nenhum curso recebeu conceito 1 e seis cursos ( 14,6 \% ) ficaram Sem Conceito ( SC ) na região Sul.
Já dos 19 cursos participantes na região Centro-Oeste ( 8,8 \% do total nacional ), oito receberam conceito 3, o conceito modal ( 42,1 \% ), como já fora destacado.
A dois cursos foram atribuídos conceito 1 ( 10,5 \% ), e a seis cursos ( 31,6 \% ), conceito 2.
O conceito 4 foi atribuído a três cursos ( 15,8 \% ).
Nessa região, nenhum curso ficou Sem Conceito ( SC ) ou recebeu conceito 5.
No Gráfico 5.1a, apresentada a distribuição cumulativa do Conceito Enade dos cursos de Licenciatura, segundo Grande Região.
Em um gráfico de distribuição cumulativa, em cada valor de abscissa ( eixo x ), o valor apresentado no eixo das ordenadas ( eixo y ) é o valor acumulado da variável.
Por exemplo, a região Norte ( linha azul ) apresenta um pouco mais de 82 \% dos cursos com conceito 2 ou abaixo deste valor.
A região Sul ( linha roxa ), por outro lado, apresenta um pouco mais de 11 \% dos cursos na mesma situação : com conceito 2 ou abaixo.
Nas regiões Norte ( linha azul ), Nordeste ( linha verde ) e Centro-Oeste ( linha preta ), todos os cursos têm conceitos 4 ou abaixo disso, já que a poligonal cumulativa alcança 100 \% para o conceito 4.
Em linhas gerais, em um gráfico de distribuição cumulativa, poligonais mais à esquerda representam Grandes Regiões com pior distribuição, e poligonais mais à direita, Grandes Regiões com melhores conceitos.
A região Norte ( linha azul ) apresenta os cursos com conceitos mais baixos ( 95,7 \% dos cursos com conceito 3 ou abaixo ), ao passo que a região Sul ( linha roxa ) apresenta a distribuição com valores maiores ( 44,4 \% dos cursos com conceito 3 ou abaixo ).
Na Tabela 5.1b, são apresentadas a quantidade e a distribuição de cursos de Bacharelado de Física participantes do Enade/2017, por faixa de conceito e Grande Região. \textbf{Enfatiza} -se, mais uma vez, que a diferença dentre os cursos tabulados neste capítulo e no capítulo 2 corresponde aos cursos Sem Conceito ( SC ), ou seja àqueles, em princípio, sem alunos concluintes que houvessem participado da prova ou àqueles que tivessem tão somente um aluno concluinte¹⁹.
Observando -se os dados da Tabela 5.1b, nota -se que, dos 57 cursos participantes, 23 ( 40,4 \% ) classificaram -se com conceito 3, o valor modal.
Esse foi também o conceito modal em quase todas as regiões : Norte ( 66,7 \% ), Sudeste ( 34,6 \% ) e Sul ( 70,0 \% ).
Nas regiões Nordeste e Centro-Oeste, a moda foi o conceito 4, atribuído a 35,7 \% e 50,0 \% dos cursos das ¹⁹ Estes cursos com somente um concluinte não foram considerados no capítulo 2. regiões, respectivamente.
O conceito 4 foi o segundo mais frequente em nível nacional ( 28,1 \%, correspondendo a 16 cursos ), e o conceito 2, o terceiro ( 17,5 \%, correspondendo a dez cursos ).
Houve, ainda, quatro cursos ( 7,0 \% ) que receberam conceito 5, e outros quatro cursos ( 7,0 \% ) que receberam conceito 1.
Dos 57 cursos de Física ( Bacharelado ), nenhum ficou Sem Conceito ( SC ).
A região Norte participou com três cursos ou 5,3 \% do total nacional.
Desses, dois cursos receberam o conceito 3, o que equivale a 66,7 \% do total regional, além de ser o conceito modal na região, como já comentado.
A um curso ( 33,3 \% ) atribuiu -se conceito 2, e nenhum curso recebeu os demais conceitos ou ficou Sem Conceito ( SC ).
A região Nordeste participou com 14 cursos ou 24,6 \% do total nacional.
Nessa região, cinco cursos ( 35,7 \% em termos regionais ) obtiveram conceito 4, o conceito modal para a região, como já comentado.
Os conceitos 1 e 2 foram atribuídos a dois cursos, cada ( 14,3 \% ).
Ainda nessa região, quatro cursos ( 28,6 \% ) foram avaliados com o conceito 3, e um curso ( 7,1 \% ), com conceito 5.
Dos 26 cursos participantes da região Sudeste ( 45,6 \% do total nacional ), o conceito 3 foi o conceito modal, como já anteriormente assinalado, tendo sido obtido por nove cursos ( 34,6 \% ).
O conceito 1 foi atribuído a dois cursos ( 7,7 \% ) e o conceito 2, a cinco cursos ( 19,2 \% ).
Receberam os conceitos 4 e 5, respectivamente, sete cursos ( 26,9 \% ) e três cursos ( 11,5 \% ).
Os dez cursos da região Sul corresponderam a 17,5 \% do total nacional.
Como já indicado, houve predominância do conceito 3, atribuído a 70,0 \% dos cursos da região ( conceito modal ), ou seja, atribuído a sete dos dez cursos participantes na região Sul.
O conceito 2 foi atribuído a um curso ( 10,0 \% ), o conceito 4, a dois cursos ( 20,0 \% ), e nenhum curso foi classificado com os demais conceitos ou ficou Sem Conceito ( SC ).
Já dos quatro cursos participantes na região Centro-Oeste ( 7,0 \% do total nacional ), dois receberam conceito 4, o conceito modal ( 50,0 \% ), como já fora destacado.
A dois cursos foram atribuídos conceitos 2 e 3 ( um curso cada ).
Nenhum curso foi classificado com os demais conceitos ou ficou Sem Conceito ( SC ).
No Gráfico 5.1b, é apresentada a distribuição cumulativa do Conceito Enade dos cursos de Bacharelado, segundo Grande Região.
Em um gráfico de distribuição cumulativa, em cada valor de abscissa ( eixo x ), o valor apresentado no eixo das ordenadas ( eixo y ) é o valor acumulado da variável.
Por exemplo, a região Norte ( linha azul ) apresenta 100 \% dos cursos com conceito 3 ou abaixo deste valor, e as regiões Sul ( linha roxa ) e Centro-Oeste ( linha preta ) apresentam 100 \% dos cursos com conceito 4 ou abaixo, pois os poligonais cumulativos atingem 100 \% no conceito 3, para a região Norte ( linha azul ), e conceito 4, nas demais regiões citadas acima.
Nota -se também que nessas três regiões citadas acima, nenhum curso obteve conceito mínimo, pois os poligonais começam a partir de 0 \% no conceito 1.
Em linhas gerais, em um gráfico de distribuição cumulativa, poligonais mais à esquerda representam Grandes Regiões com pior distribuição, e poligonais mais à direita, Grandes Regiões com melhores conceitos.
A região Norte ( linha azul ) apresenta os cursos com conceitos mais baixos, ao passo que as demais regiões do Brasil não apresentam uma hierarquia. 5.2 CONCEITOS POR CATEGORIA ADMINISTRATIVA, POR MODALIDADE DE ENSINO E GRANDE REGIÃO Na Tabela 5.2a, é apresentada a distribuição dos cursos de Licenciatura de Física participantes do Enade/2017, por Categoria Administrativa e Modalidade de ensino, de acordo com os conceitos por eles alcançados, segundo as Grandes Regiões brasileiras.
Dos 217 cursos participantes, 191 ( 88,0 \% ) eram ministrados em Instituições Públicas, e 26 ( 12,0 \% ), em Instituições Privadas.
Quanto à Modalidade de ensino, apenas 19 cursos eram a Distância ( 8,8 \% ).
De acordo com as informações que constam na Tabela 5.2a, em termos nacionais, entre Instituições Públicas e Privadas, observa -se que, dos nove cursos avaliados com conceito 5, oito eram oferecidos em IES Públicas e um, em IES Privada.
Dos 191 cursos participantes de IES Públicas, o conceito 3 foi o valor modal, atribuído a 70 cursos ( 36,6 \% ).
Dentre os demais cursos de Instituições Públicas participantes, sete cursos ( 3,7 \% ) foram avaliados com conceito 1, outros 58 cursos ( 30,4 \% ) com conceito 2, 40 cursos ( 20,9 \% ) com conceito 4 e, como já comentado, oito cursos ( 4,2 \% ) com conceito 5.
Nessa categoria, oito cursos ( 4,2 \% ) ficaram Sem Conceito ( SC ).
Na rede privada, o conceito modal também foi 3, atribuído a nove cursos, o correspondente a 34,6 \% dos 26 cursos da categoria.
Dentre os demais cursos participantes, um ( 3,8 \% ) recebeu conceito 1, e cinco ( 19,2 \% ), receberam conceito 2.
O conceito 4 foi atribuído a seis cursos ( 23,1 \% ), e o conceito 5, a um curso ( 3,8 \% ).
Nesta Categoria Administrativa, quatro cursos ( 15,4 \% ) ficaram Sem Conceito ( SC ).
Dos 198 cursos participantes da modalidade presencial, o conceito 3 foi o valor modal, atribuído a 73 cursos ( 36,9 \% ).
Dentre os demais cursos nesta modalidade de ensino, 53 cursos ( 52,6 \% ) foram avaliados com conceito 2, 45 cursos ( 22,7 \% ) obtiveram conceito 4, oito cursos ( 4,0 \% ), conceito 1, e nove cursos ( 4,0 \% ), conceito 5.
Nesta modalidade, dez cursos ( 5,1 \% ) ficaram Sem Conceito ( SC ).
Dentre os cursos a Distância, o conceito modal foi também 2, atribuído a dez cursos, o correspondente a 52,6 \% dos 19 cursos da modalidade.
Dentre os demais, seis cursos participantes ( 31,6 \% ) receberam conceito 3 e um curso ( 5,3 \% ) recebeu conceito 4.
Nesta Modalidade, dois cursos ( 10,5 \% ) ficaram Sem Conceito ( SC ), e nenhum recebeu conceito 1 ou 5.
Na análise por região, observa -se que, na região Norte, as Instituições Públicas participaram com 21 cursos ( 91,3 \% do total regional ), dos quais nenhum ficou Sem Conceito ( SC ) ou recebeu conceito 5.
O conceito modal para as Instituições Públicas na região foi 2, atribuído a 14 cursos, correspondendo a 66,7 \% dessa categoria na região.
Seis cursos foram avaliados com conceitos 1 e 3 ( três cursos, 14,3 \%, cada ), e um curso, com conceito 4 ( 4,8 \% ).
As Instituições Privadas participaram com dois cursos na região Norte ( 8,7 \% do total regional ), os quais obtiveram ambos o conceito 2, ou seja, nenhum curso recebeu os demais conceitos.
Nenhum curso ficou Sem Conceito ( SC ).
A quase totalidade dos cursos oferecidos na região Norte eram presenciais : 21 ( 91,3 \% ).
Enquanto apenas dois eram a Distância ( 8,7 \% ).
A distribuição dos cursos presenciais coincide com a das IES Públicas assim como a dos cursos a Distância coincide com a das IES Privadas nessa região.
Na região Nordeste, a rede Pública concentrou 64 dos 67 cursos participantes, o equivalente a 95,5 \% do total da região, a proporção mais \textbf{elevada} do que a registrada nas demais regiões brasileiras.
Dentre os cursos oferecidos por IES Públicas no Nordeste, 29 cursos foram avaliados com conceito 2, o conceito modal ( 45,3 \% ).
Outros 27 cursos obtiveram conceito 3 ( 42,2 \% ), dois cursos receberam conceito 1 ( 3,1 \% ), e três cursos ficaram com conceito 4 ( 4,7 \% ).
Nesta combinação de categoria e região, nenhum curso recebeu conceito 5, e três cursos ( 4,7 \% ) ficaram Sem Conceito ( SC ).
As Instituições Privadas da região Nordeste participaram com três cursos ( 4,5 \% do total da região ).
Desses, os conceitos 2, 3 e 4 foram atribuídos a um curso ( 33,3 \% ), cada.
Nenhum curso ficou Sem Conceito ( SC ).
A grande maioria dos cursos na região Nordeste também foram presenciais : 62 ( 92,5 \% ).
Para esses, a moda foi o conceito 3, com 28 cursos ( 45,2 \% ).
Os demais cursos presenciais da região foram alocados aos conceitos 1 ( dois cursos, 3,2 \% ), 2 ( 26 cursos, 41,9 \% ) e 4 ( três cursos, 4,8 \% ).
Nenhum curso presencial recebeu conceito 5, e três cursos ( 4,8 \% ) ficaram Sem Conceito ( SC ).
Os cinco cursos a Distância ( 7,5 \% ) foram avaliados com conceito 2 ( quatro cursos, 80,0 \%, conceito modal ) e com conceito 4 ( um curso, 20,0 \% ).
Nenhum curso a Distância da região foi alocado aos demais conceitos ou ficou Sem Conceito ( SC ).
Na região Sudeste, a proporção de cursos da rede Pública atingiu 88,1 \%, correspondendo a 59 dos 67 cursos participantes.
Nessa categoria e região, o conceito modal foi 3, atribuído a 25 cursos ( 42,4 \% ).
Dois cursos ( 3,4 \% ) ficaram Sem Conceito ( SC ) e nenhum recebeu conceito 1.
Os demais foram avaliados com conceito 2 ( seis cursos, 10,2 \% ), conceito 4 ( 20 cursos, 33,9 \% ) e conceito 5 ( seis cursos, 10,0 \% ).
Entre os 8 cursos ( 11,9 \% ) oferecidos por Instituições Privadas na região Sudeste, o conceito modal foi 3, atribuído a três cursos ( 37,5 \% ).
Os demais cursos receberam conceito 1, um curso ( 12,5 \% ), conceito 2, dois cursos ( 25,0 \% ) e conceito 4, um curso ( 12,5 \% ).
Um curso ficou Sem Conceito ( SC ) e nenhum curso recebeu conceito 5.
Também na região Sudeste, a maioria dos cursos foi presenciais : 59 ( 88,1 \% ).
Para esses, a moda foi o conceito 3, atribuído a 23 cursos ( 39,0 \% ).
Os demais cursos presenciais da região foram alocados aos conceitos 1 ( um curso, 1,7 \% ), 2 ( seis cursos, 10,2 \% ), 4 ( 21 cursos, 35,6 \% ) e 5 ( seis cursos, 10,2 \% ).
Somente dois cursos ( 3,4 \% ) ficaram Sem Conceito ( SC ).
Para os oito cursos a Distância ( 11,9 \% ), foram atribuídos apenas os conceitos 2 e 3, respectivamente, para dois ( 25,0 \% ) e cinco ( 62,5 \% ) cursos.
Nenhum curso a Distância da região foi alocado aos demais conceitos, e um curso ( 12,5 \% ) ficou Sem Conceito ( SC ).
As Instituições Públicas concentraram 30 dos 41 cursos participantes da região Sul, 73,2 \% do total regional, a menor das proporções desta categoria dentre as regiões.
Desses, 13 ( 43,3 \% ) obtiveram conceito 4, o conceito modal.
Nesta combinação de Categoria Administrativa e Grande Região, os conceitos 3 e 2 foram atribuídos a nove ( 30,0 \% ) e a três ( 10,0 \% ) cursos, respectivamente.
Dois cursos ( 6,7 \% ) receberam conceito 5, nenhum recebeu conceito 1, e três ( 10,0 \% ) ficaram Sem Conceito ( SC ).
As Instituições Privadas na região Sul participaram com 11 cursos, 26,8 \% dos cursos da região, a maior das proporções desta categoria nas regiões.
O conceito modal também foi 4, atribuído a quatro cursos ( 36,4 \% ).
Aos demais cursos foram atribuídos os conceitos 3 ( três cursos, 27,3 \% ) e 5 ( um curso, 9,1 \% ).
Três cursos ( 27,3 \% ) da região Sul oferecidos por IES Privadas ficaram Sem Conceito ( SC ), e nenhum curso recebeu conceito 1 ou 2.
Também na região Sul, a maioria dos cursos foram presenciais, correspondendo a 39 ( 95,1 \% ), em oposição aos dois cursos ( 4,9 \% ) a Distância.
Para os cursos presenciais, a moda foi o conceito 4, atribuído a 17 cursos ( 43,6 \% ).
Os demais cursos presenciais da região foram alocados aos conceitos 3 ( 11 cursos, 28,6 \% ) e 5 ( três cursos, 7,7 \% ), 2 ( três cursos, 7,7 \% ).
Cinco cursos presenciais ( 12,8 \% ) ficaram Sem Conceito ( SC ) e nenhum curso recebeu conceito 1.
Para os dois cursos a Distância ( 4,9 \% ), um recebeu conceito 3 e o outro ficou Sem Conceito ( SC ).
Na região Centro-Oeste, 17 dos 19 cursos participantes eram de Instituições Públicas ( 89,5 \% em termos regionais ).
Desses, seis ( 35,3 \% ) receberam conceito 2, e outros seis ( 35,3 \% ), conceito 3, os conceitos modais.
Dentre os demais, dois ( 11,8 \% ) receberam conceito 1, e três ( 17,6 \% ) receberam conceito 4.
Nenhum curso ficou Sem Conceito ( SC ) ou com conceito 5.
Quanto aos dois cursos oferecidos por Instituições Privadas na região Centro-Oeste ( 10,5 \% do total regional ), todos os cursos apresentaram conceito 3.
Nessa região, 17 cursos foram presenciais ( 89,5 \% ), e dois, a Distância ( 10,5 \% ).
Para os cursos presenciais, a moda foi o conceito 3, atribuído a oito cursos ( 47,1 \% ).
Os demais foram alocados aos conceitos 1 ( dois cursos, 11,8 \% ), 2 ( quatro cursos, 23,5 \% ) e 4 ( três cursos, 17,6 \% ).
Nenhum curso presencial ficou Sem Conceito ( SC ) ou recebeu conceito 5.
Os dois únicos cursos a Distância oferecidos na região Centro-Oeste foram avaliados com conceito 2.
No Gráfico 5.2a, é apresentada a distribuição cumulativa do Conceito Enade dos cursos de Licenciatura, segundo a Categoria Administrativa da IES.
Os cursos em IES Públicas ( linha azul ) e IES Privadas ( linha verde ) apresentam poligonais cumulativas bem semelhantes, porém nota -se que IES Públicas possuem uma porcentagem menor de cursos com conceito 1.
Na Tabela 5.2b, é apresentada a distribuição dos cursos de Bacharelado participantes do Enade/2017 de Física, por Categoria Administrativa e Modalidade de ensino, de acordo com os conceitos por eles alcançados, segundo as Grandes Regiões brasileiras.
Dos 57 cursos participantes, 53 ( 93,0 \% ) eram ministrados em Instituições Públicas, e quatro ( 7,0 \% ), em Instituições Privadas.
Quanto à Modalidade de ensino, todos os cursos eram Presenciais.
Nenhum curso, independentemente de Modalidade de ensino e Categoria Administrativa, ficou Sem Conceito ( SC ).
De acordo com as informações da Tabela 5.2b, em termos nacionais, entre Instituições Públicas e Privadas, observa -se que, dos quatro cursos avaliados com conceito 5, três eram oferecidos em IES Públicas, e um, em IES Privada.
Dos 53 cursos participantes de IES Públicas, o conceito 3 foi o valor modal, atribuído a 22 cursos ( 41,5 \% ).
Dentre os demais cursos de Instituições Públicas participantes, três cursos ( 5,7 \% ) foram avaliados com conceito 1, outros três cursos ( 5,7 \% ), com conceito 5, e 16 cursos ( 30,2 \% ), com conceito 4.
O conceito 2 foi atribuído a nove cursos ( 17,0 \% ). 1, outros três cursos ( 5,7 \% ), com conceito 5, e 16 cursos ( 30,2 \% ), com conceito 4.
O conceito 2 foi atribuído a nove cursos ( 17,0 \% ).
Na rede privada, os quatro cursos participantes ficaram \textbf{distribuídos} em quatro conceitos diferentes : 1, 2, 3 e 5.
Nenhum curso recebeu o conceito 4 ou ficou Sem Conceito ( SC ).
Conforme mencionado anteriormente, todos os cursos eram presenciais, logo, ocorre a mesma distribuição descrita a nível nacional.
Na análise por região, observa -se que, na região Norte, todos os três cursos participantes eram de Instituições Públicas ( 100 \% do total regional ), dos quais dois receberam o conceito 3 ( conceito modal ) e um recebeu o conceito 2.
Não houve participação de cursos de Instituições Privadas na região Norte.
Na região Nordeste, também não houve participação de cursos de Instituições Privadas.
Todos os 14 cursos participantes eram da rede Pública de ensino, dos quais cinco ( 35,7 \% ) foram avaliados com conceito 4, o conceito modal.
Quatro cursos obtiveram conceito 3 ( 28,6 \% ), dois cursos ( 14,3 \% ) receberam o conceito 1, dois cursos ( 14,3 \% ), o conceito 2 e um curso ( 7,1 \% ) alcançou o conceito máximo, 5.
Nenhum curso ficou Sem Conceito ( SC ).
Na região Sudeste, a proporção de cursos da rede Pública foi de 92,3 \%, correspondendo a 24 dos 26 cursos participantes.
Nesta categoria e região, o conceito modal foi 3, atribuído a nove cursos ( 37,5 \% ).
Os demais foram avaliados com conceito 4 ( sete cursos, 29,2 \% ), conceito 2 ( cinco cursos, 20,8 \% ), conceito 5 ( dois cursos, 8,3 \% ) e conceito 1 ( um curso, 4,2 \% ).
Dentre os dois cursos oferecidos por Instituições Privadas na região Sudeste, um recebeu conceito 1 e outro, conceito 5.
As Instituições Públicas concentraram oito dos dez cursos participantes da região Sul, 80,0 \% do total regional, a menor das proporções desta categoria entre as regiões.
Desses, seis obtiveram conceito 3 ( 75,0 \% ), o conceito modal, e os outros dois ( 25,0 \% ), conceito 4.
As Instituições Privadas na região Sul participaram com dois cursos, 20,0 \% dos cursos da região, a maior das proporções desta categoria nas regiões.
Os conceitos 2 e 3 foram atribuídos a um curso, cada.
Na região Centro-Oeste, todos os quatro cursos participantes eram de Instituições Públicas.
Desses, dois receberam conceito 4 ( 50,0 \% ), o conceito modal, um recebeu conceito 2 e um, também, conceito 3.
No Gráfico 5.2b, é apresentada a distribuição do Conceito Enade dos cursos de Bacharelado, segundo a Categoria Administrativa da IES.
Os cursos em IES Públicas ( linha azul ) apresentam uma poligonal mais à direita do que os em IES Privadas ( linha verde ) até o conceito 3.
Após este conceito, a configuração se inverte, e a poligonal referente aos cursos de IES Privadas passa a ficar à direita. 5.3 CONCEITOS POR ORGANIZAÇÃO ACADÊMICA E GRANDE REGIÃO Na Tabela 5.3a, encontra -se a distribuição dos conceitos atribuídos aos cursos de Licenciatura participantes do Enade/2017 na Área de Física, por Organização Acadêmica, segundo as Grandes Regiões brasileiras.
Dos 217 cursos de Física participantes, 149 eram oferecidos em Universidades, quatro em Centros Universitários, seis em Faculdades e 58 em CEFET/IFET.
Essa distribuição corresponde a, respectivamente, 68,7 \%, 1,8 \%, 2,8 \% e 26,7 \% dos cursos.
De acordo com os dados apresentados, de todos os nove cursos avaliados com conceito 5, oito eram vinculados a Universidades.
Os cursos vinculados a esse tipo de Organização Acadêmica tiveram o conceito 2 como conceito modal, atribuído a 47 cursos ( 31,5 \% ).
Os demais cursos vinculados a Universidades receberam os conceitos 1 ( quatro cursos, 2,7 \% ), 3 ( 46 cursos, 30,9 \% ), 4 ( 33 cursos, 22,1 ) e conceito 5 ( oito cursos, 5,4 \%, como já mencionado ).
Ficaram Sem Conceito ( SC ) onze cursos ( 7,4 \% ).
Dentre os cursos em Centros Universitários, os cursos ficaram \textbf{distribuídos} entre os conceitos 2 e 3, dois cursos cada.
Nenhum curso vinculado a esse tipo de Organização Acadêmica ficou Sem Conceito ( SC ) ou recebeu os demais conceitos.
Dos seis cursos mantidos por Faculdades, três ( 50,0 \% ) receberam conceito 2, o conceito modal.
Os outros três cursos oferecidos por essa categoria de IES foram alocados aos conceitos 1 e 4, recebendo dois e um curso, respectivamente.
Nenhum curso ficou Sem Conceito ( SC ) ou recebeu os conceitos 3 e 5.
Dos 58 cursos mantidos por CEFET/IFET, 31 ( 53,4 \% ) receberam conceito 3, o conceito modal.
Os conceitos 2 e 4 foram atribuídos a 11 ( 19,0 \% ) e a 12 cursos ( 20,7 \% ), respectivamente.
Dois os cursos ( 3,4 \% ) receberam conceito 1, e um curso ( 1,7 \% ) recebeu conceito 5.
Um curso ( 1,7 \% ) ficou Sem Conceito ( SC ).
Considerando -se separadamente as regiões brasileiras, verifica -se que, na região Norte, as Universidades participaram com 15 dos 23 cursos participantes ( 65,2 \% ).
Desses 15 cursos, 11 ( 73,3 \% ) receberam o conceito 2, conceito modal.
Os demais cursos receberam os conceitos 1 ( um curso, 6,7 \% ) e 3 ( três cursos, 20,0 \% ).
Nenhum curso ficou Sem Conceito ( SC ) ou recebeu conceitos 4 ou 5 nesta combinação de categoria e região.
Os Centros Universitários da região Norte não foram representados.
As Faculdades participaram com dois cursos na região Norte, que receberam o conceito 2.
Os CEFET/IFET participaram com seis cursos, 26,1 \% dos cursos da região, dos quais o conceito modal foi 2 ( três cursos, 50,0 \% ).
O conceito 1 foi atribuído a dois cursos ( 33,3 \% ), e o conceito 4, a um curso ( 16,7 \% ).
Nenhum curso ficou Sem Conceito ( SC ) ou recebeu conceitos 3 ou 5 nesta combinação de categoria e região.
Na região Nordeste, as Universidades participaram com 42 dos 67 cursos da Área de Física ( Licenciatura ) oferecidos na região, 62,7 \% dos cursos da região, dos quais dois ( 4,8 \% ) ficaram Sem Conceito ( SC ).
Nota -se que o conceito modal foi 2, atribuído a 24 cursos ( 57,1 \% ).
Obtiveram o conceito 3, 13 cursos ( 31,0 \% ), conceito 1, um curso ( 2,4 \% ) e conceito 4, dois cursos ( 4,8 \% ).
Nenhum curso oferecido por Universidades no Nordeste recebeu conceito 5.
Os Centros Universitários contaram com apenas um curso participante na região Nordeste, ao qual foi atribuído o conceito 2.
As Faculdades foram representadas por dois cursos nessa região, um recebeu conceito 1, e outro, conceito 4.
Os CEFET/IFET participaram com 22 cursos, 32,8 \% dos cursos participantes dessa região ( o maior percentual dessa organização entre as regiões ), dos quais 15 ( 68,2 \% ) receberam conceitos 3 ( conceito modal ), cinco ( 22,7 \% ), conceito 2, e um curso ( 4,5 \% ) recebeu o conceito 4.
Nenhum curso recebeu conceitos 1 ou 5, e um curso ( 4,5 \% ) ficou Sem Conceito ( SC ).
Na região Sudeste, as Universidades concentraram 46 ( 68,7 \% ) dos 67 cursos de Física ( Licenciatura ) da região.
Dentre os cursos em Universidades, nessa região, os conceitos modais foram 3 e 4, atribuídos a 16 cursos cada ( 34,8 \%, cada ).
Os demais cursos receberam os conceitos 2 ( cinco cursos, 10,9 \% ) e 5 ( seis cursos, 13,0 \% ).
Três cursos ( 6,5 \% ) ficaram Sem Conceito ( SC ) e nenhum curso recebeu conceito 1.
Os Centros Universitários participaram com dois cursos na região Sudeste, dos quais um obteve conceito 2, e outro, conceito 3.
As Faculdades também foram representadas por dois cursos na região Sudeste, dos quais um recebeu conceito 1, e outro, conceito 2.
Os CEFET/IFET participaram com 17 cursos, 25,4 \% dos cursos dessa organização na região, dos quais 11 ( 64,7 \% ) receberam o conceito modal 3.
O conceito 2 foi atribuído a um curso ( 5,9 \% ), e o conceito 4, a cinco cursos ( 29,4 \% ).
Nenhum curso ficou Sem Conceito ( SC ), nem recebeu conceitos 1 ou 5.
Dos 41 cursos da região Sul, 30 ( 73,2 \% ) eram de Universidades, para os quais o conceito modal foi 4, atribuído a 12 cursos ( 40,0 \% ).
Nesse tipo de organização, seis cursos ( 20,0 \% ) ficaram Sem Conceito ( SC ).
Os demais cursos receberam os conceitos 2 e 5 ( dois cursos, 6,7 \% cada ), e 3 ( oito cursos, 26,7 \% ).
O único curso participante de Centros Universitários da região Sul recebeu o conceito 3.
As Faculdades na região Sul não obtiveram representação.
Nessa região, foram avaliados ainda dez cursos em CEFET/IFET, que representaram 24,4 \% dos cursos participantes, e tiveram como conceito modal o conceito 4, atribuído a cinco cursos ( 50,0 \% ).
Os conceitos 2 e 5 foram atribuídos a um curso ( 10,0 \% ), cada, e o conceito 3 a três cursos ( 30,0 \% ).
Não houve curso Sem Conceito ( SC ) ou com conceito 1 nesta Organização Acadêmica da região Sul.
Na região Centro-Oeste, 16 dos 19 cursos eram de Universidades, que representaram 84,2 \% dos cursos participantes nessa região, a maior proporção que essa organização alcançou entre as regiões.
Neste tipo de organização, nenhum curso ficou Sem Conceito ( SC ) ou recebeu conceito 5.
O conceito 1 foi atribuído a dois cursos ( 12,5 \% ), o conceito 2, a cinco cursos ( 31,3 \% ), e o conceito 3, conceito modal, a seis cursos ( 37,5 \% ).
Além disso, três cursos ( 18,8 \% ) receberam o conceito 4.
Não houve representação de cursos vinculados a Centros Universitários ou Faculdades na região Centro-Oeste.
Nessa região, dos três cursos representando CEFET/IFET, dois receberam conceito 3, e um recebeu conceito 2.
No Gráfico 5.3a, é apresentada a distribuição do Conceito Enade dos cursos de Licenciatura, segundo a Organização Acadêmica da IES.
Os cursos em CEFET/IFET ( linha roxa ) e em Universidades ( linha azul ) apresentam poligonais mais à direita do que as demais, denotando uma distribuição melhor de Conceitos Enade.
No outro extremo, encontramos os cursos oferecidos em Faculdades ( linha vermelha ) e em Centros Universitários ( linha verde ) com as poligonais mais à esquerda e, consequentemente, uma distribuição de Conceitos Enade com menores valores.
Vale ressaltar que todos os cursos em Centros Universitários ( linha verde ) e em Faculdades ( linha vermelha ) estão com conceito 3, ou abaixo, e 4, ou abaixo, respectivamente.
Esse fato é notado, pois os poligonais cumulativos dessas Organizações Acadêmicas alcançam 100 \% nos conceitos 3 e 4, respectivamente.
Na Tabela 5.3b, encontra -se a distribuição dos conceitos atribuídos aos cursos de Bacharelado participantes do Enade/2017 na Área de Física, por Organização Acadêmica, segundo as Grandes Regiões brasileiras.
Dos 57 cursos de Física ( Bacharelado ) participantes, 54 eram oferecidos em Universidades, um em Centro Universitário, um em Faculdade, e um em, CEFET/IFET.
Essa distribuição corresponde a, respectivamente, 94,7 \% e 1,8 \% para cada uma das demais organizações.
Nota -se que nenhum curso ficou Sem Conceito ( SC ) nessa Habilitação.
De acordo com os dados apresentados, de todos os quatro cursos avaliados com conceito 5, todos eram vinculados a Universidades.
Os cursos vinculados a esse tipo de Organização Acadêmica tiveram o conceito 3 como conceito modal, atribuído a 23 cursos ( 42,6 \% ).
Os demais cursos vinculados a Universidades receberam os conceitos 4 ( 15 cursos, 27,8 \% ), 2 ( nove cursos, 16,7 \% ), 5 ( quatro cursos, 7,4 \%, como comentado ) e 1 ( três cursos, 5,6 \% ).
O curso vinculado a Centro Universitário recebeu conceito 2, o curso ligado à Faculdade recebeu conceito 1, e o curso em CEFET/IFET recebeu conceito 4.
Considerando -se separadamente as regiões brasileiras, verifica -se que, na região Norte, todos os três cursos participantes eram de Universidades, dos quais dois receberam conceito 3, e um, conceito 2.
Na região Nordeste, também todos os 14 cursos participantes da Área de Física ( Bacharelado ) eram de Universidades, dos quais cinco ( 35,7 \% ) receberam conceito 4, o conceito modal, quatro cursos ( 28,6 \% ) receberam o conceito 3, dois cursos ( 14,3 \% ) receberam o conceito 1, dois cursos ( 14,3 \% ), o conceito 2 ( dois cursos cada ), e um curso ( 7,1 \% ) recebeu o conceito 5.
Na região Sudeste, as Universidades concentraram 24 ( 92,3 \% ) dos 26 cursos de Física ( Bacharelado ) da região.
Dentre os cursos em Universidades, nessa região, o conceito modal foi 3, atribuído a nove cursos ( 37,5 \% ).
Os demais cursos receberam os conceitos 1 ( um curso, 4,2 \% ), 2 ( cinco cursos, 20,8 \% ), 4 ( seis cursos, 25,0 \% ) e 5 ( três cursos, 12,5 \% ).
As demais Organizações representadas na região Sudeste, com um curso cada, foram as Faculdades e CEFET/IFET, que receberam os conceitos 1 e 4, respectivamente.
Dos dez cursos da região Sul, nove ( 90,0 \% ) eram de Universidades, para os quais o conceito modal foi 3, atribuído a sete cursos ( 77,8 \% ).
Nesse tipo de organização, dois cursos receberam conceito 4 ( 22,2 \% ).
Os demais conceitos não foram atribuídos a nenhum curso.
Fora a representação das Universidades, somente os Centros Universitários participaram na região Sul, com apenas um curso, que recebeu o conceito 2.
Na região Centro-Oeste, todos os quatro cursos participantes eram de Universidades.
Nesse tipo de organização, o conceito 4 foi atribuído a dois cursos ( 50,0 \% ), conceito modal, e os conceitos 2 e 3 foram atribuídos a um curso ( 25,0 \% ), cada.
No Gráfico 5.3b, é apresentada a distribuição do Conceito Enade dos cursos de Bacharelado, segundo a Organização Acadêmica da IES.
Os cursos em Faculdades ( linha vermelha ), Centros Universitários ( linha verde ) e CEFET/IFET ( linha roxa ) obtiveram apenas uma avaliação dentro de cada Organização Acadêmica, a saber : 1, 2 e 4, respectivamente.
Os cursos em Universidades ( linha azul ) apresentam uma distribuição maior dentre os conceitos, onde 64,8 \% dos cursos obtiveram conceito 3 ou abaixo.

% matched lemmas: distribuir, elevado, enfatizar
\end{textsample}
